% Part: first-order-logic
% Chapter: introduction
% Section: sentences

\documentclass[../../../include/open-logic-section]{subfiles}

\begin{document}

\olfileid{fol}{int}{snt}

\olsection{\usetoken{P}{sentence}}

We hebben nu in grote lijnen bepaald wanneer !!{structure}s en
!!{formula}s elkaar vervullen, dus wanneer een formule ``waar is in''
een structuur. Daarvoor hadden we wel iets extra's nodig: !!a{variable}
toekenning. Maar we wilden juist !!{sentence}s definiëren. Hoe komen we
van die toekenningen af, en wat zijn !!{sentence}s precies?

Uit een eerste cursus formele logica herinner je je misschien het
verband tussen !!{variable}s en kwantoren. Na een kwantor staat altijd
!!a{variable}. Die kwantor geldt voor een bepaald deel van de
!!{sentence}; dat deel heet zijn ``bereik''. Binnen dat bereik is de
!!{variable} aan de kwantor ``gebonden''. In !!{formula}s hoeft niet
iedere !!{variable} een bijbehorende kwantor te hebben. !!{variable}s
zonder zo'n kwantor heten ``vrij'' of ``ongebonden''. We noemen precies
die !!{formula}s zonder vrije !!{variable}s voortaan !!{sentence}s.

Meestal zie je vrij snel of een voorkomen van !!a{variable} in
!!a{formula}~$!A$ vrij of gebonden is, en welke kwantor het bindt.
Misschien heb je geleerd dit te controleren door haakjes te tellen.
Voor het vervolg hebben we een precieze definitie nodig. We hebben
!!{formula}s met inductie gedefinieerd. Daarom kunnen we ook met
inductie vastleggen welke voorkomens van !!a{variable}~$x$ in
!!a{formula}~$!A$ vrij en gebonden zijn. Voor onze eenvoudige taal
werkt dat zo:
\begin{enumerate}
  \item Is $!A$ atomair, dan zijn alle voorkomens van $x$ daarin vrij.
  Het voorkomen van $x$ in~$\Atom{\Obj P}{x}$ is dus vrij.
  \item Heeft $!A$ de vorm $\lnot !B$, dan is een voorkomen van~$x$
  in~$\lnot !B$ vrij precies als het bijbehorende voorkomen van~$x$
  vrij is in~$!B$. De vrije voorkomens van variabelen in~$!B$ zijn dus
  precies de bijbehorende voorkomens in~$\lnot !B$.
  \item Heeft $!A$ de vorm $(!B \land !C)$, dan is een voorkomen van~$x$
  in~$(!B \land !C)$ vrij precies als het bijbehorende voorkomen van~$x$
  vrij is in~$!B$ of in~$!C$.
  \item Heeft $!A$ de vorm $\lexists[x][!B]$, dan is geen enkel voorkomen
  van $x$ in~$!A$ vrij. Heeft de formule de vorm $\lexists[y][!B]$ en is $y$
  een andere !!{variable} dan~$x$, dan is een voorkomen van~$x$ in
  $\lexists[y][!B]$ vrij precies als het bijbehorende voorkomen van~$x$
  vrij is in~$!B$.
\end{enumerate}

Zodra precies vastligt welke voorkomens van variabelen vrij of
gebonden zijn, is de definitie kort: !!a{sentence} is elke !!{formula}
zonder vrije voorkomens van !!{variable}s.

\end{document}
